\documentclass[a4paper,12pt]{article}

% -------------------------------------------------
% Кодировка и русский язык
% -------------------------------------------------
\usepackage[T2A]{fontenc}
\usepackage[utf8]{inputenc}
\usepackage[russian]{babel}

% -------------------------------------------------
% Математика и химические формулы
% -------------------------------------------------
\usepackage{amsmath,amssymb,mathtools}
\usepackage[version=4]{mhchem}
\usepackage{icomma}

% -------------------------------------------------
% Шрифт
% -------------------------------------------------
\usepackage{paratype}
\renewcommand{\familydefault}{\sfdefault}

% -------------------------------------------------
% Вёрстка
% -------------------------------------------------
\usepackage[
    a4paper,
    left=20mm,
    right=20mm,
    top=24mm,
    bottom=22mm,
    headheight=25pt,
    headsep=8mm
]{geometry}

\usepackage{setspace}
\usepackage{parskip}
\usepackage{enumitem}
\usepackage{tabularx,booktabs}
\usepackage{array}
\usepackage{xcolor}
\usepackage{fancyhdr}
\usepackage{hyperref}

% -------------------------------------------------
% Графика
% -------------------------------------------------
\usepackage{tikz}
\usetikzlibrary{arrows.meta,positioning,calc}

% -------------------------------------------------
% Цветовая палитра
% -------------------------------------------------
\definecolor{accent}{HTML}{008080}
\definecolor{darkaccent}{HTML}{004D4D}
\definecolor{textgray}{HTML}{333333}
\definecolor{softgray}{HTML}{F1F4F4}
\definecolor{linegray}{HTML}{AAB8B8}

% -------------------------------------------------
% Общие настройки
% -------------------------------------------------
\onehalfspacing
\setlength{\parindent}{0pt}
\setlength{\parskip}{0.55em}
\setlength{\emergencystretch}{3em}
\raggedbottom

\hypersetup{
    unicode=true,
    colorlinks=true,
    linkcolor=darkaccent,
    urlcolor=accent,
    pdftitle={Алканы: способы получения},
    pdfauthor={Лёвин Артём Александрович}
}

\setlist[itemize]{
    leftmargin=*,
    itemsep=0.25em,
    topsep=0.3em
}

\setlist[enumerate]{
    leftmargin=*,
    itemsep=0.65em,
    topsep=0.4em
}

\newcolumntype{Y}{>{\raggedright\arraybackslash}X}

% -------------------------------------------------
% Колонтитулы
% -------------------------------------------------
\pagestyle{fancy}
\fancyhf{}
\fancyhead[L]{\small\color{textgray}Алканы: способы получения}
\fancyhead[R]{\small\color{textgray}Ксения Клыкова}
\fancyfoot[C]{\small\color{textgray}\thepage}
\renewcommand{\headrulewidth}{0.4pt}
\renewcommand{\headrule}{%
    \hbox to\headwidth{%
        \color{linegray}%
        \leaders\hrule height \headrulewidth\hfill
    }%
}

% -------------------------------------------------
% Пользовательские команды
% -------------------------------------------------
\newcommand{\sectiontitle}[1]{%
    \vspace{0.9em}%
    {\Large\bfseries\color{darkaccent}#1\par}%
    \vspace{0.2em}%
    {\color{accent}\rule{\linewidth}{0.8pt}}%
    \vspace{0.45em}%
}

\newcommand{\methodtitle}[2]{%
    \vspace{0.9em}%
    {\large\bfseries\color{darkaccent}#1. #2\par}%
    \vspace{0.25em}%
}

\newcommand{\important}[1]{%
    \noindent
    \begin{tikzpicture}
        \node[
            fill=softgray,
            draw=linegray,
            rounded corners=2pt,
            inner sep=7pt,
            text width=0.94\linewidth,
            align=left
        ] {\textbf{\color{darkaccent}Важно.} #1};
    \end{tikzpicture}%
}

\begin{document}

% =================================================
% ТИТУЛЬНЫЙ ЛИСТ
% =================================================
\begin{titlepage}
    \thispagestyle{empty}
    \centering

    \vspace*{2.3cm}

    {\Large\bfseries\color{accent}ЧЕК-ЛИСТ\par}

    \vspace{0.65cm}

    {\Huge\bfseries\color{darkaccent}
        Алканы: способы получения
        \par
    }

    \vspace{0.5cm}

    {\Large\color{textgray}
        Промышленные и лабораторные методы
        \par
    }

    \vspace{2.2cm}

    {\large
        Лёвин Артём Александрович
        \par
    }

    \vspace{0.25cm}

    {\large\bfseries
        эксклюзивно для Ксении Клыковой
        \par
    }

    \vspace{0.35cm}

    {\normalsize
        10 класс
        \par
    }

    \vfill

    {\large\color{textgray}\today\par}

    \vspace{1.4cm}

    \begin{tikzpicture}[x=1cm,y=1cm]
        \draw[
            color=accent,
            line width=1.1pt
        ]
        (-7,0)
        .. controls (-4.8,1.05) and (-2.6,-0.8) ..
        (0,0)
        .. controls (2.8,0.9) and (4.8,-0.9) ..
        (7,0);

        \draw[
            color=linegray,
            line width=0.6pt
        ]
        (-6.3,-0.28)
        .. controls (-3.5,0.45) and (-1.6,-0.55) ..
        (1,-0.15)
        .. controls (3.1,0.2) and (4.5,-0.45) ..
        (6.2,-0.25);
    \end{tikzpicture}

\end{titlepage}

% =================================================
% ОСНОВНОЙ ЧЕК-ЛИСТ
% =================================================
\sectiontitle{1. Главная идея занятия}

\textbf{Алканы} --- предельные ациклические углеводороды, в молекулах которых
между атомами углерода присутствуют только одинарные связи.

Общая формула алканов:
\[
    \mathrm{C_nH_{2n+2}}, \qquad n\geq 1.
\]

При выборе способа получения алкана прежде всего необходимо определить,
как должен измениться его углеродный скелет:

\begin{itemize}
    \item сохранить число атомов углерода;
    \item удвоить исходный углеродный фрагмент;
    \item уменьшить цепь на один атом углерода;
    \item получить метан из простых веществ или карбида;
    \item выделить алканы из природной смеси.
\end{itemize}

\begin{center}
\begin{tikzpicture}[
    >=Latex,
    node distance=9mm and 10mm,
    main/.style={
        draw=darkaccent,
        fill=softgray,
        rounded corners=3pt,
        line width=0.9pt,
        text width=4.1cm,
        align=center,
        inner sep=8pt,
        font=\bfseries
    },
    branch/.style={
        draw=linegray,
        rounded corners=2pt,
        text width=4.6cm,
        align=center,
        inner sep=6pt
    },
    arrow/.style={
        ->,
        line width=0.8pt,
        color=accent
    }
]
    \node[main] (center) {Получение алканов};

    \node[branch, above left=of center] (oil)
        {Выделение из смеси\\перегонка нефти};

    \node[branch, left=of center] (same)
        {Скелет сохраняется\\гидрирование};

    \node[branch, below left=of center] (minus)
        {Цепь уменьшается на один \(\mathrm C\)\\реакция Дюма};

    \node[branch, above right=of center] (methane)
        {Получение метана\\прямой синтез, гидролиз карбидов};

    \node[branch, right=of center] (double)
        {Фрагмент удваивается\\реакция Вюрца};

    \node[branch, below right=of center] (kolbe)
        {Фрагмент удваивается\\электролиз Кольбе};

    \draw[arrow] (center) -- (oil);
    \draw[arrow] (center) -- (same);
    \draw[arrow] (center) -- (minus);
    \draw[arrow] (center) -- (methane);
    \draw[arrow] (center) -- (double);
    \draw[arrow] (center) -- (kolbe);
\end{tikzpicture}
\end{center}

\sectiontitle{2. Основные способы получения алканов}

\methodtitle{1}{Природный газ и нефть}

Природный газ и нефть являются важнейшими промышленными источниками
углеводородов. Природный газ содержит преимущественно метан, а нефть
представляет собой сложную смесь углеводородов различного строения.

Для разделения компонентов нефти применяют
\textbf{фракционную перегонку}, или \textbf{ректификацию}.

\textbf{Принцип метода:} вещества разделяются благодаря различию
температур кипения.

При нагревании:

\begin{enumerate}
    \item более летучие компоненты переходят в пар;
    \item пары поднимаются по ректификационной колонне;
    \item на разных уровнях колонны пары конденсируются;
    \item образуются фракции с различными интервалами температур кипения.
\end{enumerate}

\important{Перегонка является физическим способом разделения смеси, а не
химической реакцией. Нефтяные фракции обычно представляют собой смеси
веществ, а не индивидуальные чистые алканы.}

В лабораторном примере смесь этанола и воды можно частично разделить,
поскольку при нормальном давлении этанол кипит приблизительно при
\(78{,}4^\circ\mathrm C\), а вода --- при \(100^\circ\mathrm C\).
Для более глубокого разделения применяют многократное испарение и
конденсацию --- ректификацию.

\methodtitle{2}{Прямой синтез метана}

Метан можно получить непосредственным взаимодействием углерода с водородом:

\[
    \ce{C + 2H2 -> CH4}.
\]

Реакция проводится при повышенной температуре, давлении и в присутствии
катализатора.

\textbf{Особенность метода:} прямой синтез из простых веществ используют
прежде всего для получения метана.

\methodtitle{3}{Каталитическое гидрирование непредельных углеводородов}

\textbf{Гидрирование} --- реакция присоединения водорода.

Алкены содержат двойную связь и имеют общую формулу
\(\mathrm{C_nH_{2n}}\). При присоединении одной молекулы водорода двойная
связь превращается в одинарную, а алкен становится алканом:

\[
    \ce{C_{n}H_{2n} + H2 -> C_{n}H_{2n+2}}.
\]

Типичные катализаторы:

\[
    \ce{Ni},\qquad \ce{Pd},\qquad \ce{Pt}.
\]

Пример:

\[
    \ce{CH2=CH2 + H2 ->[Ni] CH3-CH3}.
\]

Этен превращается в этан.

Другой пример:

\[
    \ce{CH3-CH=CH-CH3 + H2 ->[Pt] CH3-CH2-CH2-CH3}.
\]

Бутен превращается в бутан.

\textbf{Алгоритм гидрирования:}

\begin{enumerate}
    \item найдите кратную связь;
    \item добавьте одну молекулу \(\ce{H2}\) на одну двойную связь;
    \item замените двойную связь на одинарную;
    \item добавьте по одному атому водорода к каждому из двух атомов углерода;
    \item проверьте четырёхвалентность углерода.
\end{enumerate}

\important{При гидрировании углеродный скелет не изменяется. Меняются только
тип связи и число атомов водорода.}

\methodtitle{4}{Реакция Вюрца}

Реакция Вюрца --- взаимодействие галогеналканов с металлическим натрием.

Общая схема:

\[
    \ce{2R-X + 2Na -> R-R + 2NaX},
\]

где \(R\) --- углеводородный фрагмент, а \(X\) --- атом галогена.

Реакцию проводят в сухом эфире.

Пример получения этана:

\[
    \ce{2CH3Cl + 2Na -> CH3-CH3 + 2NaCl}.
\]

Пример получения бутана:

\[
    \ce{2CH3-CH2Cl + 2Na ->
        CH3-CH2-CH2-CH3 + 2NaCl}.
\]

\textbf{Смысл реакции:}

\begin{itemize}
    \item натрий связывается с атомами галогена;
    \item два одинаковых углеводородных фрагмента соединяются друг с другом;
    \item число атомов углерода в цепи удваивается.
\end{itemize}

Если исходный галогеналкан содержит \(k\) атомов углерода, полученный
алкан содержит \(2k\) атомов углерода.

\important{Для получения одного чистого продукта предпочтительно использовать
одинаковые галогеналканы. При взаимодействии двух разных галогеналканов
образуется смесь нескольких алканов.}

\clearpage

\sectiontitle{3. Дополнительные способы получения}

\methodtitle{5}{Гидролиз карбидов}

\textbf{Гидролиз} --- взаимодействие вещества с водой, сопровождающееся
его химическим превращением.

\textbf{Карбиды} --- бинарные соединения углерода с металлами или некоторыми
менее электроотрицательными элементами.

При гидролизе карбида алюминия образуется метан:

\[
    \ce{Al4C3 + 12H2O -> 4Al(OH)3 + 3CH4 ^}.
\]

При гидролизе карбида бериллия также образуется метан:

\[
    \ce{Be2C + 4H2O -> 2Be(OH)2 + CH4 ^}.
\]

В этих карбидах углерод условно рассматривают как частицу со степенью
окисления \(-4\). При присоединении атомов водорода образуется молекула
\(\ce{CH4}\).

\important{Не все карбиды при гидролизе образуют метан. Например, карбид
кальция \(\ce{CaC2}\) реагирует с водой с образованием ацетилена
\(\ce{C2H2}\), а не алкана.}

\methodtitle{6}{Декарбоксилирование солей карбоновых кислот. Реакция Дюма}

При нагревании натриевой соли карбоновой кислоты с гидроксидом натрия
в присутствии оксида кальция образуется алкан:

\[
    \ce{RCOONa + NaOH -> RH + Na2CO3},
    \qquad \ce{CaO},\ \Delta.
\]

Смесь \(\ce{NaOH}\) и \(\ce{CaO}\) называют \textbf{натронной известью}.

Пример получения метана:

\[
    \ce{CH3COONa + NaOH -> CH4 + Na2CO3},
    \qquad \ce{CaO},\ \Delta.
\]

Пример получения этана:

\[
    \ce{CH3CH2COONa + NaOH ->
        CH3CH3 + Na2CO3},
    \qquad \ce{CaO},\ \Delta.
\]

\textbf{Декарбоксилирование} --- удаление карбоксильной группы с потерей
одного атома углерода.

Если соль содержит \(m\) атомов углерода, то полученный алкан содержит
\(m-1\) атом углерода.

\important{В реакции Дюма углеродная цепь уменьшается ровно на один атом
углерода.}

\methodtitle{7}{Электролиз солей карбоновых кислот. Реакция Кольбе}

\textbf{Электролиз} --- совокупность окислительно-восстановительных реакций,
протекающих на электродах при прохождении электрического тока через раствор
или расплав электролита.

Общая схема электролиза водного раствора натриевой соли одноосновной
карбоновой кислоты:

\[
    \ce{2RCOONa + 2H2O ->
        R-R + 2CO2 ^ + H2 ^ + 2NaOH}.
\]

Пример для ацетата натрия:

\[
    \ce{2CH3COONa + 2H2O ->
        CH3-CH3 + 2CO2 ^ + H2 ^ + 2NaOH}.
\]

Образуется этан.

Пример для пропионата натрия:

\[
    \ce{2CH3CH2COONa + 2H2O ->
        CH3CH2-CH2CH3 + 2CO2 ^ + H2 ^ + 2NaOH}.
\]

Образуется бутан.

\textbf{Главная закономерность:} два одинаковых углеводородных фрагмента
\(R\) соединяются между собой.

Если исходная соль содержит \(m\) атомов углерода с учётом атома углерода
карбоксильной группы, то образующийся алкан содержит

\[
    2(m-1)
\]

атомов углерода.

\important{При электролизе одной соли образуются преимущественно симметричные
алканы с чётным числом атомов углерода.}

\sectiontitle{4. Сводная таблица}

\renewcommand{\arraystretch}{1.35}

\begin{table}[ht]
\centering
\caption{Изменение углеродного скелета при получении алканов}
\small
\begin{tabularx}{\linewidth}{@{}p{3.2cm}Y Y p{2.7cm}@{}}
\toprule
\textbf{Способ} &
\textbf{Исходное вещество} &
\textbf{Изменение скелета} &
\textbf{Ключевое условие}
\\
\midrule

Перегонка нефти &
Смесь углеводородов &
Химическое строение не изменяется; смесь разделяется на фракции &
Разные температуры кипения
\\

Прямой синтез &
Углерод и водород &
Получается метан &
Температура, давление, катализатор
\\

Гидрирование &
Алкен или другой непредельный углеводород &
Число атомов углерода сохраняется &
\(\ce{H2}\), \(\ce{Ni}\), \(\ce{Pd}\) или \(\ce{Pt}\)
\\

Реакция Вюрца &
Галогеналкан &
Два фрагмента соединяются; число атомов углерода удваивается &
\(\ce{Na}\), сухой эфир
\\

Гидролиз карбидов &
Метаниды металлов &
Получается метан &
Вода
\\

Реакция Дюма &
Натриевая соль карбоновой кислоты &
Цепь уменьшается на один атом углерода &
\(\ce{NaOH}\), \(\ce{CaO}\), нагревание
\\

Реакция Кольбе &
Соль карбоновой кислоты &
Два фрагмента \(R\) соединяются &
Вода, электрический ток
\\

\bottomrule
\end{tabularx}
\end{table}

\sectiontitle{5. Универсальный алгоритм выбора реакции}

\begin{enumerate}
    \item Запишите формулу алкана, который необходимо получить.
    \item Определите число атомов углерода в продукте.
    \item Сравните углеродный скелет продукта с возможным исходным веществом.
    \item Выберите требуемое изменение:
    \begin{itemize}
        \item сохранить скелет --- гидрирование;
        \item удвоить фрагмент --- Вюрц или Кольбе;
        \item уменьшить цепь на один атом углерода --- Дюма;
        \item получить метан --- прямой синтез или гидролиз метанидов;
        \item разделить природную смесь --- перегонка.
    \end{itemize}
    \item Запишите реагенты и условия реакции.
    \item Расставьте коэффициенты.
    \item Проверьте число атомов каждого элемента в левой и правой частях.
    \item Убедитесь, что каждый атом углерода образует четыре связи.
\end{enumerate}

\sectiontitle{6. Типичные ошибки}

\begin{itemize}
    \item Называть перегонку химической реакцией.
    \item Считать, что при перегонке нефти получают только индивидуальные
    чистые вещества.
    \item Забывать водород или катализатор при гидрировании.
    \item Сохранять двойную связь после присоединения \(\ce{H2}\).
    \item Нарушать четырёхвалентность углерода.
    \item Забывать коэффициент \(2\) перед галогеналканом и натрием
    в реакции Вюрца.
    \item Пытаться получить чистый несимметричный алкан реакцией Вюрца
    из двух разных галогеналканов.
    \item Считать, что любой карбид при гидролизе образует метан.
    \item Забывать, что в реакции Дюма цепь уменьшается на один атом углерода.
    \item Не записывать побочные продукты электролиза Кольбе:
    \(\ce{CO2}\), \(\ce{H2}\) и \(\ce{NaOH}\).
\end{itemize}

\sectiontitle{7. Что необходимо знать в первую очередь}

Для базового школьного уровня необходимо уверенно распознавать и применять:

\begin{enumerate}
    \item выделение углеводородов из природного газа и нефти;
    \item прямой синтез метана;
    \item каталитическое гидрирование;
    \item реакцию Вюрца.
\end{enumerate}

Гидролиз карбидов, реакция Дюма и электролиз Кольбе особенно полезны
в заданиях повышенной сложности и в цепочках превращений.

% =================================================
% ТРЕНИРОВОЧНЫЙ БЛОК
% =================================================
\clearpage
\sectiontitle{Тренировочный блок}

\textbf{Упражнения 1--3 --- средний уровень.}

\begin{enumerate}
    \item Запишите уравнение гидрирования пропена. Назовите полученный алкан.

    \item Запишите уравнение реакции Вюрца для хлорметана.
    Назовите органический продукт.

    \item Расставьте коэффициенты:
    \[
        \ce{Al4C3 + H2O -> Al(OH)3 + CH4}.
    \]

    \item Запишите структурную формулу продукта полного гидрирования
    бутена-2. Укажите один подходящий катализатор.

    \item Запишите уравнение реакции Вюрца для 1-бромпропана.
    Назовите полученный алкан.

    \item Какой алкан образуется при нагревании бутаноата натрия
    \(\ce{CH3CH2CH2COONa}\) с натронной известью?
    Запишите уравнение реакции.

    \item Запишите уравнение электролиза водного раствора пропионата натрия
    \(\ce{CH3CH2COONa}\). Укажите все продукты реакции.

    \item Получите этан тремя различными способами:
    \begin{itemize}
        \item гидрированием непредельного углеводорода;
        \item реакцией Вюрца;
        \item реакцией Дюма.
    \end{itemize}
    Запишите уравнения реакций.

    \item Определите, какие утверждения являются ошибочными. Исправьте их.
    \begin{enumerate}[label=\alph*)]
        \item При перегонке нефть разделяется на индивидуальные чистые алканы.
        \item При взаимодействии хлорметана и хлорэтана с натрием
        образуется только пропан.
        \item При гидролизе карбида кальция образуется метан.
    \end{enumerate}

    \item Предложите три независимых способа получения бутана:
    \begin{itemize}
        \item гидрированием подходящего алкена;
        \item реакцией Вюрца;
        \item электролизом Кольбе.
    \end{itemize}
    Для каждого способа выберите исходное вещество, запишите уравнение
    и укажите необходимые условия.
\end{enumerate}

% =================================================
% ОТВЕТЫ
% =================================================
\clearpage
\sectiontitle{Ответы к тренировочному блоку}

\renewcommand{\arraystretch}{1.4}

\begin{table}[ht]
\centering
\caption{Краткие ответы}
\small
\begin{tabularx}{\linewidth}{@{}p{1.1cm}Y@{}}
\toprule
\textbf{№} & \textbf{Ответ} \\
\midrule

1 &
\(\ce{CH3CH=CH2 + H2 ->[Ni] CH3CH2CH3}\).
Продукт --- пропан.
\\

2 &
\(\ce{2CH3Cl + 2Na -> CH3CH3 + 2NaCl}\).
Продукт --- этан.
\\

3 &
\(\ce{Al4C3 + 12H2O -> 4Al(OH)3 + 3CH4}\).
\\

4 &
\(\ce{CH3CH=CHCH3 + H2 ->[Pt] CH3CH2CH2CH3}\).
Продукт --- бутан.
\\

5 &
\(\ce{2CH3CH2CH2Br + 2Na ->
CH3CH2CH2CH2CH2CH3 + 2NaBr}\).
Продукт --- гексан.
\\

6 &
\(\ce{CH3CH2CH2COONa + NaOH ->
CH3CH2CH3 + Na2CO3}\),
\(\ce{CaO}\), нагревание.
Продукт --- пропан.
\\

7 &
\(\ce{2CH3CH2COONa + 2H2O ->
CH3CH2CH2CH3 + 2CO2 + H2 + 2NaOH}\).
Продукт --- бутан.
\\

8 &
Гидрирование:
\(\ce{CH2=CH2 + H2 ->[Ni] CH3CH3}\).

Реакция Вюрца:
\(\ce{2CH3Cl + 2Na -> CH3CH3 + 2NaCl}\).

Реакция Дюма:
\(\ce{CH3CH2COONa + NaOH ->
CH3CH3 + Na2CO3}\),
\(\ce{CaO}\), нагревание.
\\

9 &
Все три утверждения ошибочны.

а) При перегонке нефти получают фракции, каждая из которых обычно является
смесью углеводородов.

б) Смесь хлорметана и хлорэтана в реакции Вюрца даёт смесь этана,
пропана и бутана.

в) \(\ce{CaC2 + 2H2O -> C2H2 + Ca(OH)2}\);
образуется ацетилен.
\\

10 &
Гидрирование:
\(\ce{CH2=CHCH2CH3 + H2 ->[Ni] CH3CH2CH2CH3}\).

Реакция Вюрца:
\(\ce{2CH3CH2Cl + 2Na ->
CH3CH2CH2CH3 + 2NaCl}\), сухой эфир.

Электролиз Кольбе:
\(\ce{2CH3CH2COONa + 2H2O ->
CH3CH2CH2CH3 + 2CO2 + H2 + 2NaOH}\),
электрический ток.
\\

\bottomrule
\end{tabularx}
\end{table}

\vfill

\begin{center}
    {\color{darkaccent}\bfseries
        Главная цель --- не механически запомнить все реакции,\\
        а научиться выбирать способ по изменению углеродного скелета.
    }
\end{center}

\end{document}
